% PAGES: 322-324

The no--arbitrage values of European call and put options with the same strike and
maturity satisfy a model--independent relationship called the Put--Call parity. The
intuition behind the Put--Call parity is that a long call and a short put position
in options with the same strike and maturity is the same as a long position in a
forward contract on the underlying asset of the options with the same expiration
date as the maturity of the options and with delivery price equal to the strike of
the options.

More formally, if $C(t)$ and $P(t)$ are the values at time $t$ of a European call
and put option, respectively, with maturity $T$ and strike $K$, on the same
underlying asset with spot price $S(t)$, paying dividends continuously at the rate
$q$, the Put--Call parity states that
\begin{equation}
C(t) \;-\; P(t) \;=\; S(t)e^{-q(T-t)} \;-\; Ke^{-r(T-t)}.
\end{equation}
A proof of (10.92) based on the Law of One Price can be found in section 1.9 of
Stefanica \cite{36}.

\medskip
\noindent\textbf{The Black--Scholes formulas}

The Black--Scholes formulas give the price of plain vanilla European call and put
options under the assumptions that the price of the underlying asset has lognormal
distribution with volatility $\sigma$, the asset pays dividends continuously at the
rate $q$, and the risk--free interest rate is constant and equal to $r$. The
Black--Scholes values $C_{BS}(S,t)$ and $P_{BS}(S,t)$ of a call option and of a put
option, respectively, with strike $K$ and maturity $T$ are given by
\begin{align}
C(S,t) &= Se^{-q(T-t)}N(d_1) \;-\; Ke^{-r(T-t)}N(d_2); \\
P(S,t) &= Ke^{-r(T-t)}N(-d_2) \;-\; Se^{-q(T-t)}N(-d_1),
\end{align}
where
\begin{align}
d_1 &= \frac{\ln\left(\frac{S}{K}\right) +
\left(r-q+\frac{\sigma^2}{2}\right)(T-t)}{\sigma\sqrt{T-t}}; \\
d_2 &= d_1 - \sigma\sqrt{T-t} \;=\; \frac{\ln\left(\frac{S}{K}\right) +
\left(r-q-\frac{\sigma^2}{2}\right)(T-t)}{\sigma\sqrt{T-t}},
\end{align}
and $N(z)$ denotes the cumulative distribution of the standard normal variable,
i.e.,
\[
N(z) \;=\; \frac{1}{\sqrt{2\pi}}\int_{-\infty}^z e^{-\frac{x^2}{2}}\,dx.
\]

\medskip
\noindent\textbf{Options Greeks}

The Greeks are the sensitivities of the options prices with respect to various
parameters, e.g., with respect to the price of the underlying asset or with respect
to the volatility of the underlying asset, and are important for hedging purposes.

From the Black--Scholes formulas (10.93) and (10.94) for options on assets paying
continuous dividends, the following closed formulas for the Greeks of European plain
vanilla call and put options can be derived:
\begin{align}
\Delta(C) &= \pd{C}{S} \;=\; e^{-q(T-t)}N(d_1); \\
\Delta(P) &= \pd{P}{S} \;=\; -e^{-q(T-t)}N(-d_1); \\
\Gamma(C) &= \pdd{C}{S} \;=\; \frac{e^{-q(T-t)}}{S\sigma\sqrt{2\pi(T-t)}}
e^{-\frac{d_1^2}{2}}; \\
\Gamma(P) &= \pdd{P}{S} \;=\; \Gamma(C); \\
\text{vega}(C) &= \pd{C}{\sigma} \;=\; Se^{-q(T-t)}\sqrt{\frac{T-t}{2\pi}}
e^{-\frac{d_1^2}{2}}; \\
\text{vega}(P) &= \pd{P}{\sigma} \;=\; \text{vega}(C).
\end{align}

\section{Eigenvalues of symmetric matrices}

In this section, we provide an elegant proof of the fact that any eigenvalue of a
symmetric matrix with real entries is a real number; cf. Theorem 5.1 and Theorem
10.8.

To do so, we introduce an extension of the inner product (5.3), also called the
Euclidean inner product, to vectors with complex entries.

\begin{definition}
Let $u = (u_i)_{i=1:n}$ and $v = (v_i)_{i=1:n}$ be column vectors of size $n$ with
entries complex numbers, i.e., with $u_i,v_i \in \C$, for $i = 1:n$. The complex
Euclidean inner product of $u$ and $v$ is\footnote{Note that, if the entries of $u$
and $v$ are real numbers, i.e., if $u_i,v_i \in \R$ for all $i = 1:n$, then the
complex Euclidean inner product is equal to the inner product given by (5.3):
\[
(u,v)_{\C} \;=\; \sum_{i=1}^n u_i\overline{v_i} \;=\; \sum_{i=1}^n u_iv_i \;=\;
(u,v).
\]}
\begin{equation}
(u,v)_{\C} \;=\; u_1\overline{v_1} + u_2\overline{v_2} + \ldots + u_n\overline{v_n}
\;=\; \sum_{i=1}^n u_i\overline{v_i},
\end{equation}
where $\overline{v_i}$ denotes the complex conjugate of $v_i$, for $i = 1:n$.
\end{definition}

Note that
\begin{equation}
(u,v)_{\C} \;=\; v^*u,
\end{equation}
where, by definition, $v^* = (\overline{v_1} \ \overline{v_2} \ \ldots \
\overline{v_n}) = (\overline{v_i})_{i=1:n}$ is the row vector whose entries are the
complex conjugates of the entries of the vector $v = (v_i)_{i=1:n}$.

The following properties of the Euclidean complex inner product follow from
definition (10.103):
\begin{align}
(u,v)_{\C} &= \overline{(v,u)_{\C}}, \quad \forall\, u,v \in \C^n; \\
(cu,v)_{\C} &= c(u,v)_{\C}, \quad \forall\, u,v \in \C^n,\ c \in \C; \\
(u,cv)_{\C} &= \overline{c}(u,v)_{\C}, \quad \forall\, u,v \in \C^n,\ c \in \C.
\end{align}

\begin{definition}
The Euclidean norm of a vector $v = (v_i)_{i=1:n} \in \C^n$ is
\begin{equation}
||v||_{\C} \;=\; \sqrt{\sum_{i=1}^n ||v_i||^2} \;=\; \sqrt{(v,v)_{\C}} \;=\;
\sqrt{v^*v}.
\end{equation}
\end{definition}

\begin{definition}
Let $A$ be an $m \times n$ matrix with complex entries. The hermitian of the matrix
$A$ is the $n \times m$ matrix $A^*$ given by
\[
A^*(j,k) \;=\; \overline{A(k,j)}, \quad \forall\, k = 1:n,\ j = 1:m.
\]
\end{definition}

\begin{lemma}
The hermitian of a matrix with real entries is the transpose of the matrix.
\end{lemma}

\begin{proof}
Let $A$ be an $m \times n$ matrix with real entries. Then, $\overline{A(j,k)} =
A(j,k)$, and therefore
\begin{equation}
A^*(k,j) \;=\; \overline{A(j,k)} \;=\; A(j,k) \;=\; A^t(k,j), \quad \forall\, j =
1:m,\ k = 1:n.
\end{equation}
From (10.109), we conclude that $A^* = A^t$.
\end{proof}

The properties below are similar to those from Lemma 1.1 and Lemma 1.2 for the
transpose of a matrix, and can be proved similarly:

For any $m \times n$ matrix $A$,
\begin{equation}
(A^*)^* \;=\; A.
\end{equation}

For any $m \times n$ matrix $A$, $n \times p$ matrix $B$, and column vector $v$ of
size $n$,
\begin{align}
(Av)^* &= v^*A^*; \\
(AB)^* &= B^*A^*.
\end{align}

\begin{lemma}
Let $A$ be a square matrix of size $n$ with complex entries, and let $u$ and $v$ be
column vectors of size $n$ with complex entries. Then,
\begin{align}
(Au,v)_{\C} &= (u,A^*v)_{\C}; \\
(u,Av)_{\C} &= (A^*u,v)_{\C}.
\end{align}
\end{lemma}

\begin{proof}
Recall from (10.111) that $(Av)^* = v^*A^*$. Since $(A^*)^* = A$, see (10.110), we
find that $(A^*v)^* = v^*(A^*)^* = v^*A$. Then, using (10.104), we obtain that
\begin{align*}
(Au,v)_{\C} &= v^*Au \;=\; (v^*A)u \;=\; (A^*v)^*u \;=\; (u,A^*v)_{\C}; \\
(u,Av)_{\C} &= (Av)^*u \;=\; v^*A^*u \;=\; v^*(A^*u) \;=\; (A^*u,v)_{\C}.
\end{align*}
\end{proof}

We can now prove Theorem 5.1 from Section 5.1:

\begin{theorem}
Any eigenvalue of a symmetric matrix with real entries is a real number.
\end{theorem}

\begin{proof}
Let $A$ be a square matrix of size $n$ with real entries, and assume that $A$ is
symmetric, i.e., $A^t = A$. Let $\lambda \in \C$ and $v \in \C^n$ be an eigenvalue
and a corresponding eigenvector of $A$. Then, $Av = \lambda v$, and, using (10.106)
and (10.108), we find that
\begin{equation}
(Av,v)_{\C} \;=\; (\lambda v,v)_{\C} \;=\; \lambda(v,v)_{\C} \;=\;
\lambda||v||_{\C}^2.
\end{equation}
% Proof continues onto PDF page 325 (folio 309), the first page of the next file
% in this chunk -- the \begin{proof} above is intentionally left open here.
